% GENERATED FILE. Edit canonical lessons, then run npm run generate:books.
% canonical-source-hash: fnv1a64:74a91bbd4d8924cf
% canonical-lessons: FR-C441-banc, FR-C441-lampadaire, FR-C441-reverbere, FR-C441-kiosque, FR-C441-pave

\chapter{Dans la rue}
\label{ch:fr-dans-la-rue}
\hlchaptermodality{441}

\begin{chapteropening}
\textbf{By the end of this chapter:} \emph{I can describe things in a street: the bench, the street lamp, the old lamp post, the kiosk and the cobblestones.}

Complete the last lesson of chapter 441: I can describe things in a street: the bench, the street lamp, the old lamp post, the kiosk and the cobblestones.
\end{chapteropening}

\section[le banc]{le banc — the bench}

\label{lesson:FR-C441-banc}



\begin{quote}
Before the new one: say the French for the police station, then the French for the barracks.
\end{quote}

\subsection*{You'll want to know: le banc}
\textbf{le banc} — \textquotedblleft{}the bench\textquotedblright{}.

\begin{quote}
Deux vieux messieurs lisent sur un banc.
\end{quote}

\begin{quote}
\emph{Two old men are reading on a bench.}
\end{quote}

\emph{Un banc de poissons} — a shoal of fish.

\begin{cousinweb}
\emph{banc} comes from a Germanic word for 'bench' (EtymDB), and English \emph{bench} is a cousin.
\end{cousinweb}

\begin{grammarlens}[title={What you've built}]
One more word for telling, explaining and describing at B1.
\end{grammarlens}

\subsection*{Your turn}
\begin{itemize}
\raggedright
  \item \emph{Say it:} \emph{le banc}
  \item \emph{Say it:} \emph{le banc}, once more
  \item \emph{Say it:} a sentence of your own with \emph{le banc}
  \item \emph{Recall:} say the French for the police station, then the French for the barracks, then say \emph{le banc} again
\end{itemize}

\begin{tcolorbox}[breakable,colback=teal!4,colframe=teal!35!black,title={Before you move on}]
What does le banc mean? (\textquotedblleft{}The bench\textquotedblright{}.) Say it once more.
\end{tcolorbox}
\section[le lampadaire]{le lampadaire — the street lamp; the floor lamp}

\label{lesson:FR-C441-lampadaire}



\begin{quote}
Before the new one: say the French for the barracks, then the French for the bench.
\end{quote}

\subsection*{You'll want to know: le lampadaire}
\textbf{le lampadaire} — \textquotedblleft{}the street lamp; the floor lamp\textquotedblright{}.

\begin{quote}
Les lampadaires s'allument à six heures.
\end{quote}

\begin{quote}
\emph{The street lamps come on at six o'clock.}
\end{quote}

At home, \emph{un lampadaire} is a standing lamp.

\begin{cousinweb}
\emph{lampadaire} is a learned borrowing from Medieval Latin \emph{lampadārium} 'candlestick', from \emph{lampas}, from Greek \emph{lampás} 'torch' (Lewis \& Short). EtymDB does not list it.
\end{cousinweb}

\begin{grammarlens}[title={What you've built}]
One more word for telling, explaining and describing at B1.
\end{grammarlens}

\subsection*{Your turn}
\begin{itemize}
\raggedright
  \item \emph{Say it:} \emph{le lampadaire}
  \item \emph{Say it:} \emph{le lampadaire}, once more
  \item \emph{Say it:} a sentence of your own with \emph{le lampadaire}
  \item \emph{Recall:} say the French for the barracks, then the French for the bench, then say \emph{le lampadaire} again
\end{itemize}

\begin{tcolorbox}[breakable,colback=teal!4,colframe=teal!35!black,title={Before you move on}]
What does le lampadaire mean? (\textquotedblleft{}The street lamp; the floor lamp\textquotedblright{}.) Say it once more.
\end{tcolorbox}
\section[le réverbère]{le réverbère — the lamp post (old-style)}

\label{lesson:FR-C441-reverbere}



\begin{quote}
Before the new one: say the French for the bench, then the French for the street lamp.
\end{quote}

\subsection*{You'll want to know: le réverbère}
\textbf{le réverbère} — \textquotedblleft{}the lamp post (old-style)\textquotedblright{}.

\begin{quote}
Dans la vieille rue, il reste quelques réverbères.
\end{quote}

\begin{quote}
\emph{In the old street a few old lamp posts are still standing.}
\end{quote}

Older word, often for old-fashioned lamps.

\begin{cousinweb}
\emph{réverbère} is built on \emph{réverbérer} 'to reflect', a learned borrowing from Latin \emph{reverberāre} 'to beat back', from \emph{verber} 'whip' (Lewis \& Short): the old lamps had a mirror that threw back the light.
\end{cousinweb}

\begin{grammarlens}[title={What you've built}]
One more word for telling, explaining and describing at B1.
\end{grammarlens}

\subsection*{Your turn}
\begin{itemize}
\raggedright
  \item \emph{Say it:} \emph{le réverbère}
  \item \emph{Say it:} \emph{le réverbère}, once more
  \item \emph{Say it:} a sentence of your own with \emph{le réverbère}
  \item \emph{Recall:} say the French for the bench, then the French for the street lamp, then say \emph{le réverbère} again
\end{itemize}

\begin{tcolorbox}[breakable,colback=teal!4,colframe=teal!35!black,title={Before you move on}]
What does le réverbère mean? (\textquotedblleft{}The lamp post (old-style)\textquotedblright{}.) Say it once more.
\end{tcolorbox}
\section[le kiosque]{le kiosque — the kiosk, the newsstand}

\label{lesson:FR-C441-kiosque}



\begin{quote}
Before the new one: say the French for the street lamp, then the French for the lamp post.
\end{quote}

\subsection*{You'll want to know: le kiosque}
\textbf{le kiosque} — \textquotedblleft{}the kiosk, the newsstand\textquotedblright{}.

\begin{quote}
Il achète le journal au kiosque du coin.
\end{quote}

\begin{quote}
\emph{He buys the paper at the kiosk on the corner.}
\end{quote}

\emph{Un kiosque à musique} — a bandstand.

\begin{cousinweb}
\emph{kiosque} was borrowed from Turkish \emph{köşk} 'pavilion', which came from Persian \emph{kušk} 'palace, portico' (EtymDB).
\end{cousinweb}

\begin{grammarlens}[title={What you've built}]
One more word for telling, explaining and describing at B1.
\end{grammarlens}

\subsection*{Your turn}
\begin{itemize}
\raggedright
  \item \emph{Say it:} \emph{le kiosque}
  \item \emph{Say it:} \emph{le kiosque}, once more
  \item \emph{Say it:} a sentence of your own with \emph{le kiosque}
  \item \emph{Recall:} say the French for the street lamp, then the French for the lamp post, then say \emph{le kiosque} again
\end{itemize}

\begin{tcolorbox}[breakable,colback=teal!4,colframe=teal!35!black,title={Before you move on}]
What does le kiosque mean? (\textquotedblleft{}The kiosk, the newsstand\textquotedblright{}.) Say it once more.
\end{tcolorbox}
\section[le pavé]{le pavé — the cobblestone; the paving}

\label{lesson:FR-C441-pave}



\begin{quote}
Before the new one: say the French for the lamp post, then the French for the kiosk.
\end{quote}

\subsection*{You'll want to know: le pavé}
\textbf{le pavé} — \textquotedblleft{}the cobblestone; the paving\textquotedblright{}.

\begin{quote}
Les vélos roulent mal sur les pavés.
\end{quote}

\begin{quote}
\emph{Bikes do not ride well on the cobbles.}
\end{quote}

\emph{Un pavé} is also a very thick book or a thick steak.

\begin{cousinweb}
\emph{pavé} is the past participle of \emph{paver} 'to pave' (EtymDB), which goes back to Latin \emph{pavīre} 'to beat the earth flat' (Lewis \& Short).
\end{cousinweb}

\begin{grammarlens}[title={What you've built}]
One more word for telling, explaining and describing at B1.
\end{grammarlens}

\subsection*{Your turn}
\begin{itemize}
\raggedright
  \item \emph{Say it:} \emph{le pavé}
  \item \emph{Say it:} \emph{le pavé}, once more
  \item \emph{Say it:} a sentence of your own with \emph{le pavé}
  \item \emph{Recall:} say the French for the lamp post, then the French for the kiosk, then say \emph{le pavé} again
\end{itemize}

\begin{tcolorbox}[breakable,colback=teal!4,colframe=teal!35!black,title={Before you move on}]
What does le pavé mean? (\textquotedblleft{}The cobblestone; the paving\textquotedblright{}.) Say it once more.
\end{tcolorbox}
